\section{Особенности реализации}

Программа реализована на языке C++. Поля числа хранит структура
\mintinline{cpp}{BFloat16}. Преобразования выполняет отдельный класс-конвертер.
Взаимодействие с пользователем обеспечивает класс \mintinline{cpp}{CLI}.
Структура хранит поля формата отдельно, что делает этапы преобразования
наглядными. Её размер в
памяти не обязан быть равен двум байтам, поскольку задача структуры ---
представить логическое устройство числа, а не оптимизировать его физическую
упаковку.

\codeListing[cpp][firstline=7,lastline=11]
    {../include/bfloat16.h}
    {Структура представления bfloat16}
    {\label{lst:structure}}

Метод \mintinline{cpp}{encode} выделяет знак числа и отдельно обрабатывает NaN,
бесконечность и ноль. Для конечного ненулевого значения строятся двоичные записи
целой и дробной частей. Метод \mintinline{cpp}{normalize} находит первую значащую единицу,
определяет истинную экспоненту и формирует последовательность битов, из которой
будет получена мантисса.

Если целая часть ненулевая, экспонента равна номеру старшего разряда целой части.
При нулевой целой части выполняется поиск первой единицы в дробной части; её
позиция определяет отрицательную экспоненту. Ведущая единица исключается из
последовательности мантиссы, поскольку для нормализованных значений она является
неявной.

\codeListing[cpp][firstline=181,lastline=201,fontsize=\scriptsize,baselinestretch=0.9]
    {../src/bfloat16.cpp}
    {Нормализация двоичного представления}
    {\label{lst:normalize}}

Метод \mintinline{cpp}{roundMantissa} оставляет семь старших битов мантиссы и анализирует
следующий разряд. Если он равен единице, сохранённое значение увеличивается на
единицу. Переполнение семиразрядной мантиссы приводит к её обнулению и переносу
единицы в экспоненту. Такой способ соответствует выбранному в программе
округлению по первому отбрасываемому биту.

Метод \mintinline{cpp}{decode} выполняет обратное преобразование. Для нормализованного
числа применяется формула~\eqref{eq:normal}; ноль, субнормальные значения,
бесконечность и NaN обрабатываются отдельными ветвями. При этом декодирование
субнормальных значений поддерживается полностью, а при прямом кодировании
значения, экспонента которых не помещается в нормализованный диапазон, в текущей
учебной реализации округляются к нулю.

Класс \mintinline{cpp}{CLI} отвечает только за взаимодействие с пользователем. Он выводит
меню, проверяет ввод и вызывает методы конвертера. Математические операции в
классе интерфейса не выполняются.

